\section{Overview}

This chapter describes the experimental methodology for examining AI persuasiveness
in a turn-based debate format. The primary analysis focuses on \textbf{AI-to-human
persuasion}: whether debating against a GPT-4o mini opponent configured with one of
three personalities (firm, balanced, open-minded) produces different belief or
confidence outcomes relative to a human-human control condition. A secondary analysis
investigates confidence change as a co-primary outcome. Exploratory analyses examine
reflection quality as a mechanism and argument quality as an alternative predictor.
This study is treated as a \textbf{pilot investigation}; the primary aim is to
estimate effect sizes and establish infrastructure for future full-scale deployment
rather than to draw confirmatory conclusions.

A key structural asymmetry exists in the design: human participants complete
between-round reflections (paraphrase and acknowledgement), whereas AI agents do not.
Reflection quality is therefore analyzable only on the human side. Analysis of
human-to-AI persuasion is deferred to future work.


\section{Study Design}

The study uses a controlled between-conditions design with repeated measures over
debate rounds. Participants are assigned to one of four conditions: three AI opponent
conditions (firm, balanced, open-minded) or a human-human control. Each debate
follows an identical round structure with arguments submitted in alternating turns
and belief and confidence checks collected after each round.


\section{Participants and Recruitment}

Participants were recruited according to institutional guidelines and provided with
an informed consent statement describing the debate task and data usage. Inclusion
required completion of the full debate protocol, including per-round belief and
confidence checks and the post-debate survey. The final analytic sample comprised
$n = 90$ participants across four conditions: human-human ($n = 21$), firm AI
($n = 28$), balanced AI ($n = 20$), and open-minded AI ($n = 21$).


\section{Conditions and Randomization}

Participants were assigned to conditions with an operational goal of maintaining
approximate balance across the three AI personalities and the human-human control.
Assignment was monitored through an administrative matchmaking workflow and adjusted
during data collection to keep condition counts as balanced as feasible. Opponent
identity was concealed during the debate to prevent source bias. Topic selection
and initial stance were self-selected by participants; some imbalance across
conditions on these dimensions may remain and is addressed by including topic
fixed effects in the primary models.


\section{Procedure}

\begin{enumerate}
    \item Participant selects a debate topic and indicates initial stance and
          confidence (0--100).
    \item Participant completes pre-debate survey recording openness to persuasion:
          \textit{firm on my stance}, \textit{convinced of my stance}, or
          \textit{open to changing my mind}.
    \item Debate proceeds in alternating argument rounds (500-character limit per
          turn). After each round, the participant records updated belief value
          (0--100) and confidence (0--100).
    \item Between rounds, the participant completes a reflection task comprising
          a free-text paraphrase of the opponent's most recent argument and a
          free-text acknowledgement of whether they concede the opponent's point.
    \item After the final round, the participant completes a post-debate survey
          including opponent attribution and confidence in that attribution.
\end{enumerate}

\pagebreak
\section{Measures}

\subsection{Primary Outcome: Belief Change}

Participant belief is captured on a continuous 0--100 scale, recorded before the
debate (Round 0) and after each subsequent round. Belief change is computed as
$\Delta \text{belief} = \text{belief}_{\text{post}} - \text{belief}_{\text{pre}}$,
where pre-debate belief is taken from Round 0 and post-debate belief from the
participant's final round response.

\subsection{Secondary Outcome: Confidence Change}

Confidence is measured on a 0--100 scale at the first post-argument round
(Round 1, because confidence was not collected at pre-debate (Round 0); \texttt{confidence\_pre})
and at the participant's final round (\texttt{confidence\_post}).
Confidence change is computed as
$\Delta \text{confidence} = \text{confidence}_{\text{post}} - \text{confidence}_{\text{pre}}$.
Baseline confidence at Round 1 shows substantial variation (M = 79.3, SD = 18.4),
but with ceiling effects \(20\% at 100\), so confidence-change models are
interpreted with this constrained variance in mind.

\subsection{Reflection Quality}

Between-round reflections were scored on an overall quality scale
using GPT-4o mini as the primary rater, drawing on dimensions of
faithfulness, specificity, engagement, clarity, and charity
collapsed into a single composite score. To validate GPT scoring,
two independent human raters scored a stratified random subsample
of $n = 30$ reflections using the same rubric. Convergent validity
was assessed using Spearman $\rho$ between the GPT score and the
mean of the two human ratings; inter-rater reliability between
human raters was assessed using two-way random effects ICC$(2,1)$
(single-rater reliability) and ICC$(2,k)$ (average-rater
reliability), implemented via mean-squares decomposition. A
combined ICC was additionally computed by treating GPT as a third
rater alongside the two humans. Rows with missing values on any
rater were excluded prior to ICC computation. Reflection quality
scores were standardized ($z$-scored) prior to inclusion as a
model covariate.

\subsection{Argument Quality}

Persuader argument quality was scored as a single overall quality
rating using GPT-4o mini, drawing on dimensions of evidence
quality, logical coherence, responsiveness, nuance, and clarity.
Validation followed an identical procedure to reflection quality:
two independent human raters scored a stratified subsample of
$n = 30$ arguments; Spearman $\rho$ was computed between GPT
scores and the mean human rating; two-way random effects
ICC$(2,1)$ and ICC$(2,k)$ were computed for human-only and
GPT$+$human matrices respectively, with listwise deletion of
incomplete rows. Argument quality scores were standardized
($z$-scored) prior to inclusion as an exploratory covariate.

\subsection{Pre-Debate Openness}

Participants completed a single-item pre-debate survey selecting one of three
openness options. This variable is treated as a baseline covariate
(\texttt{initial\_stance\_strength\_z}) capturing resistance to persuasion and
is included in confirmatory models M4 and M5.

\subsection{Acknowledgement}

The free-text acknowledgement field records whether the participant explicitly
conceded validity in the opponent's argument. This was binarized
(\texttt{acknowledgement\_binary}) and included as an exploratory covariate in
a secondary belief change model (\texttt{belief\_ack}).

\subsection{Opponent Attribution}

The post-debate survey records opponent perception (human, AI, or unsure),
perception confidence (1--5), suspicion timing (round range), detection cues
(multi-select), and AI awareness effect (more persuasive, less persuasive, no
difference). These variables are available for exploratory analysis but are
not part of the primary confirmatory model sequence.


\section{Data Quality and Exclusion Criteria}

Debates were excluded if incomplete, missing per-round belief checks, or
containing inconsistent round sequencing. The descriptive sample is $n = 90$,
while complete-case model analyses use $n = 89$. No implausible
response time exclusions were applied; all exclusions are reported
transparently.

\pagebreak
\section{Statistical Analysis}

\subsection{Pre-Model Checks}

Prior to confirmatory model fitting, four diagnostic checks were
conducted.
First, zero-mass on \texttt{belief\_change} was audited
overall (14.4\%) and by condition (human-human: 23.8\%; firm: 10.7\%;
balanced: 15.0\%; open-minded: 9.5\%).
Second, a Shapiro-Wilk test
confirmed approximate normality of \texttt{belief\_change}
($W = 0.980$, $p = 0.190$), supporting OLS use.
Third, a one-way ANOVA confirmed no significant between-condition
difference in \texttt{belief\_change} prior to covariate adjustment ($F = 0.488$,
$p = 0.691$); AI-subcondition ANOVA yielded $\eta^2 = 0.001$
($H = 0.606$, $p = 0.739$, Kruskal-Wallis), confirming pooling is
defensible.
Fourth, topic fixed effects were justified by an adjusted
$R^2$ increment of 0.070 over condition-only models, confirming topics
are a meaningful confounder. An \texttt{nrounds} imbalance was noted:
the balanced condition averaged 5.80 rounds vs.\ 4.19 for human-human
($SD = 2.44$ vs.\ $0.51$), motivating its inclusion as a covariate in
M5.

\subsection{Rationale for Pooling AI Conditions}

The three AI personality conditions (firm, balanced, open-minded) are pooled
into a single binary indicator (\texttt{group\_2lvl}: AI vs.\ human-human)
for the primary confirmatory analysis. This decision is motivated by three
considerations. First, the pilot sample ($n = 89$) yields per-condition
cell sizes of $n = 20$--$28$, which are insufficient to reliably estimate
separate personality effects; pooling recovers adequate power for the
higher-priority human-vs-AI contrast. Second, the primary research question
concerns whether AI debate \textit{per se} produces different outcomes than
human debate, independent of personality configuration; personality-level
differences are a secondary question addressed in Section~\ref{sec:part10}.
Third, pre-model ANOVA across AI subconditions confirmed negligible
between-personality differences in belief change ($\eta^2 = 0.001$,
Kruskal-Wallis $p = 0.739$), making pooling theoretically defensible.
 Condition-level trajectories are examined in the
four-level analysis (Section~\ref{sec:part10}), which retains the full
condition structure.


\subsection{Primary Confirmatory Model Sequence}

Belief change is modelled using OLS regression with HC3 heteroskedasticity-robust
standard errors. The focal predictor is a binary group indicator (AI pooled
vs.\ human-human), with human-human as the reference category. Seven models are
estimated (four confirmatory, three exploratory/sensitivity):

\begin{itemize}
    \item \textbf{M1}: Treatment effect only — unadjusted estimate of AI
          vs.\ human-human belief change.
    \item \textbf{M2}: Treatment $+$ reflection quality ($z$-scored) —
          adjusting for participant engagement depth.
     \item \textbf{M3} (exploratory): Treatment $\times$ reflection quality
          interaction — testing whether reflection quality moderates the
          effect of AI debate on belief change.
    \item \textbf{M4}: Treatment $+$ reflection quality $+$ initial stance
          strength ($z$-scored) $+$ topic fixed effects — the primary
          fully-adjusted confirmatory model.
    \item \textbf{M4b} (exploratory): Treatment $+$ reflection quality $+$
          persuader argument quality ($z$-scored) — an omitted-variable
          check testing whether argument quality accounts for variance
          beyond reflection quality alone, without topic or stance controls.
    \item \textbf{M4c} (exploratory): Treatment $+$ persuader argument
          quality $+$ initial stance strength $+$ topic fixed effects —
          substituting argument quality for reflection quality to isolate
          the independent predictive contribution of persuader argument
          quality.
    \item \textbf{M5}: Full model — treatment $+$ reflection quality $+$
          argument quality $+$ initial stance strength $+$ topic fixed
          effects $+$ number of rounds $+$ first-player indicator, serving
          as a kitchen-sink sensitivity check for coefficient stability.
\end{itemize}

Given the pilot model sample ($n = 89$),
observed power for the M1 treatment effect was 0.20 at the observed
Cohen's $d = 0.286$; a total $n \approx 350$ (175 per group) would
be required for 80\% power at this effect size.
Effect sizes are reported as standardized mean differences (Cohen's $d$) with
95\% confidence intervals. Given the pilot sample size ($n = 89$), observed
power and required sample sizes for 80\% power are reported to contextualize
null findings.

\subsection{Secondary Analysis: Confidence Change}

Confidence change models mirror the primary belief change sequence and are
estimated using the same OLS/HC3 framework. Two confirmatory and two
exploratory models are specified:

\begin{itemize}
    \item \textbf{conf-M1} (confirmatory): Unadjusted treatment effect on
          confidence change — AI vs.\ human-human, no covariates.
    \item \textbf{conf-M2} (confirmatory): Treatment $+$ reflection quality
          ($z$-scored) — adjusting for engagement depth, mirroring the
          primary M2 specification.
    \item \textbf{conf-M3} (exploratory): Treatment $\times$ reflection
            quality interaction on confidence change — included for
            symmetry with the primary M3 specification; the interaction
            term is not expected to be interpretable given the
            variance-constrained nature of confidence change.
    \item \textbf{belief-ack} (exploratory): Belief change regressed on
        treatment $+$ reflection quality $+$ acknowledgement binary
        (\texttt{acknowledgement\_binary}) — testing whether explicit
        concession behaviour predicts belief change over and above
        reflection quality ($n = 89$, no missingness on this variable).

\end{itemize}

\texttt{belief\_ack} is estimated on the available complete-case sample
($n = 89$ in the current export); results are interpreted as exploratory
and are not subject to multiplicity correction.

\subsection{Robustness Checks}

All confirmatory models are subjected to:
\begin{itemize}
    \item \textbf{Outlier trimming}: Re-estimation after excluding the bottom
          and top 1\% of belief or confidence change values.
    \item \textbf{Leave-one-out (LOO)}: Re-estimation of M2 excluding each
          condition in turn.
    \item \textbf{VIF diagnostics}: Variance inflation factors computed for
          all multi-predictor models to detect multicollinearity.
    \item \textbf{Shapiro-Wilk}: Residual normality assessed for all models.
\end{itemize}

\subsection{Exploratory Analyses}

Exploratory analyses are not part of primary inference and are reported
as hypothesis-generating. Three categories of exploratory analysis are
conducted.

\subsubsection{Reflection Quality as Mediator}

A qualitative subsample analysis examines the thematic content of
between-round reflections to contextualise the quantitative reflection
quality scores. The subsample is stratified by reflection quality tercile
(low, medium, high) and AI condition, capped at two observations per
debate to avoid within-participant duplication, yielding approximately
$n = 24$ reflections. Reflections are coded for themes including
argument acknowledgement, counter-argumentation style, and evidence
of perspective-taking. This analysis is intended to validate whether
GPT-assigned reflection quality scores correspond to substantively
meaningful differences in engagement, rather than surface-level
textual features.

\subsubsection{Mechanism Analysis: Pushback and Confidence}

If co-primary confidence effects are observed, a mechanism analysis
tests whether critical engagement --- operationalised as pushback
frequency (\texttt{has\_pushback}) --- predicts confidence change.
Pushback is coded from the qualitative subsample (the same
$n \approx 24$ reflections used in the thematic analysis) as the
binary presence of disagreement markers (\textit{but},
\textit{however}, \textit{disagree}, and synonyms). These qualitative
codes are merged with participant-level outcome measures
(\texttt{confidence\_change}, \texttt{belief\_change}) to form the
mechanism dataset.

Pearson correlations between \texttt{has\_pushback} and
\texttt{confidence\_change} are estimated separately within the AI
group and the human-human group. Fisher-$z$ transformation is used
to compute 95\% confidence intervals around each correlation. A
negative correlation within the AI group would suggest that
participants who engaged more critically gained \textit{less}
confidence, contradicting a defensive confidence hypothesis; a
positive correlation would suggest active pushback reinforces
perceived debate success.

Jackknife resampling (leave-one-out) is applied to the AI-group
correlation to assess sensitivity to individual influential
observations. At each iteration, one participant is removed and the
correlation is re-estimated on the remaining subsample; if the
sign remains consistent across all iterations, the correlation is
classified as robust. Iterations where pushback or confidence
change variance collapses to zero are excluded.

Prior to mechanism analysis, qualitative-sample representativeness is
audited by comparing condition proportions in the subsample against
full-population condition proportions. In the current run, the
subsample overrepresents human-human debates and underrepresents AI
conditions; mechanism results are therefore interpreted as
hypothesis-generating rather than population-representative.


\section{Reporting}

Results are reported with effect sizes and 95\% confidence intervals rather
than $p$-values alone. All models report coefficient estimates, robust standard
errors, and sample sizes. Exploratory analyses are clearly labelled and
interpreted with appropriate caution. All statistical code, intermediate tables,
and export files are retained to facilitate reproducibility. Cross-chapter
headline numbers are synchronized to Appendix Table~\ref{tab:headline_numbers}.
